\paragraph{Short harmonic polynomials on every column lattice.}
Let \(O=O_E\), \(e=[E:\mathbb Q]=2f\), \(R_E=|\Delta_E|^{1/e}\),
and let \(F\) denote the coefficient field of the chosen deep
cover (in particular it contains the labelled images of \(E\)).
Write \(\Phi\) for the chosen CM type. In the following calculation
\(\Lambda\subset E^n\) is the \emph{column} lattice, self-dual for
\(v^*u\). It may vary through the entire allowed genus. The period
lattices of the abelian varieties have rank \(r=p+q\) over \(O\).

\begin{proposition}[A common short family for all embeddings]
\label{bfam:harm}
Assume \(n\ge p+q\). At each \(\tau\in\Phi\) prescribe two
partitions of lengths at most \(p,q\), respectively. Let \(d_P\)
be the sum of their sizes. Consider row highest weight polynomials
in the two matrices of \(p,q\) rows and \(n\) columns at each
\(\tau\), annihilated by all mixed contractions
\(\sum_b\partial_{X_{ib}}\partial_{Y_{jb}}\) at each type.
This space over \(F\) is spanned by a finite set of polynomials
in the integral algebra generated by the linear forms in the rows
with coefficients \(\tau(u)\) in the first block and
\(\tau(\bar v)\) in the second, for \(u,v\in\Lambda\).
Put \(D=(2p+1)d_P,\ h=en,\ T=\max(1,hD\sqrt e\,nR_E)\). Every
polynomial \(P\) in the set satisfies, at every embedding of \(F\),
\[
       \|P^\sigma\|_1\le (r\sqrt n\,T^{2p+1})^{d_P},
 \qquad \#\{P\}\le (D+1)^{hr}.
\]
The norm is the sum of absolute coefficients in the original
column coordinates. For the stage parameters \(L\ge1\),
\(\log(fr)\le L,\ \log R_E\le3L,\ \log S\le20L\), where
\(S=d_P+frn/2\), this gives
\(\log^+\|P^\sigma\|_1\ll rL d_P\), with an absolute constant.
\end{proposition}
\begin{proof}
Use the real Euclidean structure with norm squared
\(\operatorname{Tr}_{E/\mathbb Q}(u^*u)\). Its dual lattice is
\(\mathfrak D_E^{-1}\Lambda\) by self-duality and the trace
description of the codifferent. For \(0\ne v\) in that dual,
\(v^*v\in\mathfrak D_E^{-2}\) since conjugation preserves the
different. Its absolute field norm is at least
\(|\Delta_E|^{-2}\) by inclusion of principal fractional ideals.
All its conjugates are positive. Their arithmetic-geometric mean
inequality gives \(\|v\|_{\rm tr}\ge\sqrt e/R_E\). Euclidean
transference in dimension \(h\) now gives \(h\) independent
\(b_i\in\Lambda\) of trace norm at most \(\sqrt e\,nR_E\).
They form a real basis of the product of the \(f\) complex
column spaces. \cite[Theorem~2.1]{banaszczyk}

Independently choose \(p+q\) vectors \(u_i,w_j\) on the grid
\(\sum_{\ell=1}^h a_\ell b_\ell,\ 0\le a_\ell\le D\),
with integral \(a_\ell\). Write \(U\) for the matrix of the \(u_i\)
as columns, \(G=U^*U\), and set
\[
       v_j=\det(G)w_j-U\,\operatorname{adj}(G)U^*w_j .
\]
The coefficients used to combine the lattice vectors are integral
by their Hermitian pairings. Also \(U^*v_j=0\) by the adjugate
identity. At any embedding the input vectors have Euclidean norm
at most \(T\). For independent columns the displayed operator is
\(\det G\) times an orthogonal projection. By continuity its norm
is at most \(T^{2p}\) even at singular tuples, by Hadamard's
inequality. Hence \(\|\tau v_j\|\le T^{2p+1}\) at every embedding.

At \(\tau\), contract each of the first rows with the ordered list
\(\tau u_i\), and each of the second rows with \(\tau\bar v_j\),
using the ordinary bilinear dot product. Take the determinants of
the first \(i\) rows and columns in each resulting square matrix,
to powers the successive differences of the respective partitions
(use the reversed column order in the standard opposite block).
Multiply over \(i\) and \(\tau\). The determinants give the
specified row character. Every mixed derivative contracts a pair
of mutually orthogonal coefficient vectors and is zero. A linear
form has coefficient norm at most \(\sqrt n\,T^{2p+1}\).
A determinant of order \(i\) costs at most an additional \(i!\le
r^i\), giving the asserted bound. Its coefficients lie in the
stated integral algebra, also for a nonfree projective \(\Lambda\).

To see spanning, a complex linear functional annihilating the
grid polynomials gives a polynomial in the \(hr\) real grid
coordinates of degree at most \(D\): conjugation is real-linear
and each projected vector has degree \(2p+1\). Successive
one-variable interpolation makes it zero for all real choices.
The \(b_i\) allow independent choices at the \(f\) archimedean
places. In particular at each place use the first \(p\) and last
\(q\) columns of any unitary matrix, in the orders of the
prescribed minors. Here \(G=1\) and the projection fixes the
last columns. The obtained products form the orbit of the
product of the standard opposite flag minors under \(U(n)^f\).
Kashiwara and Vergne identify this product as a highest weight
vector and its column representation as the entire harmonic
row highest weight space for the given partitions, with multiplicity
one \cite[unitary case, \S6, Proposition~6.1 and Theorem~6.3]{kashiwara-vergne}.
Taking products gives the assertion at the \(f\) places. Descent
of matrix rank gives spanning over \(F\), since the equations
and grid vectors are defined there. Finally \(n\le2S,\ d_P\le S\)
give \(\log T\ll L\) if \(d_P>0\). For \(d_P=0\) take just \(1\).
\end{proof}
The two sets of rows here are labelled by
\(\tau,\bar\tau\), with the dual rows reversed when passing to
the ordered flag weight. Thus the integral algebra is exactly
the one generated by lattice directions in the harmonic jet of
Theorem~\ref{intjet:semi}. Relabelling the letters by an embedding
permutes the column coordinates and preserves the coefficient
estimate. For each fixed auxiliary integer, lattice and level
tuple the theta sum is linear in \(P\): the jet is the evaluation
of the constant-coefficient operator \(P(\partial)\) and the
translations and finite sums are independent of \(P\). Consequently
the short family has exactly the same image span as all such
harmonic polynomials for these data. This costs no determinant
or denominator from a change of basis.
\src{src/Development/AbelianVarietyTheory/Proofs/CuspFormEvaluations/ExistsShortHarmonicFamilyUnimodular.lean}{96}{The integral algebra and its simultaneous coefficient bounds};
\src{src/Development/AbelianVarietyTheory/Proofs/DivisorProofs/PolynomialCoefficientEmbeddingNormLeOfL1Bound.lean}{399}{Use at the actual deep coefficient field}.

\paragraph{From a half maximum to the polarized period metric.}
The transfer from the lattice in Proposition~\ref{bergvol:global}
to the periods in Proposition~\ref{intjet:arch} must account for
the host fractional ideal. Write the actual period lattice as
\(\Pi=\mathfrak a L'\), where \(L'\subset E^r\) is self-dual
for the row Hermitian form \(h_0\). Write \(\xi\) for its
imaginary polarization coefficient and \(x=|\tau\xi|\),
independent of \(\tau\). The row polarization is
\(\operatorname{Tr}_{E/\mathbb Q}(-\xi h_0)\), with
\((\xi)\mathfrak a\bar{\mathfrak a}\mathfrak D_E=O\).

\begin{lemma}[Comparison of the two minima]\label{bfam:minimum}
Let \({\rm maj}_\tau\) denote the positive majorant of \(h_0\)
at an orthonormal flag. Assume that
\[
        \sum_{\tau\in\Phi}{\rm maj}_\tau(u)\ge\varepsilon
                 \qquad(0\ne u\in L').
\]
Then the shortest nonzero polarized period
has squared length at least \(2\varepsilon/R_E^2\).
\end{lemma}
\begin{proof}
Taking ideal norms gives \(N(\mathfrak a)^{-1/f}=xR_E\).
In the product of \(f\) complex planes, \(\mathfrak a^{-1}\)
has covolume \(2^{-f}\sqrt{|\Delta_E|}/N(\mathfrak a)\).
The product of discs of squared radius \(xR_E^2\) has volume
strictly greater than \(2^{2f}\) times this covolume, as
\(\pi>2\). Minkowski gives \(0\ne b\in\mathfrak a^{-1}\)
with \(|\tau b|^2\le xR_E^2\) simultaneously. For \(v\in\Pi\)
nonzero, apply the hypothesis to \(bv\in L'\).
The polarized squared norm of \(v\) is
\(2x\sum_\tau{\rm maj}_\tau(v)\): in the orthogonal positive
and negative coordinates one conjugates the positive block in the
period convention of Lemma~\ref{dg:gaussian}, and the two terms
in the trace give \(2x\) on both
blocks. This also follows from the imaginary part of the Hermitian
form, which is the displayed trace polarization. The estimate
follows on summing
\({\rm maj}_\tau(bv)=|\tau b|^2{\rm maj}_\tau(v)\).
\end{proof}
\src{src/Development/AbelianVarietyTheory/Proofs/DivisorProofs/FirstMinAtHalfMax.lean}{82}{The inverse ideal and the distinction between the two lattices}.

\begin{corollary}[Simultaneous estimate in the normalized cover metric]
\label{bfam:metric}
Use the framed uniformization of the deep cover at each good embedding
and the self-dual row lattice model above. Assume the analytic cusp
reading and the half-maximum condition of the first hypothesis of
Proposition~\ref{bergvol:global} for the theta sections in question.
Use the parameters of Corollary~\ref{bergvol:size} and
Proposition~\ref{intjet:arch}, with \(S=d_P+frn/2\) and
\(\log S\le20L\). The sections read from the polynomials of
Proposition~\ref{bfam:harm}, multiplied by the reflex-ideal
scalar \(s\) of Theorem~\ref{intjet:trace}, have logarithmic
pointwise supremum at most \(C(frL)^2S\) in the normalized cover
metric, for an absolute \(C\). The same bound holds for the full
cover Petersson norm and every positive-mass component Petersson
norm. Multiplication by the common integer of
Theorem~\ref{intjet:semi} adds only its logarithm.
\end{corollary}
\begin{proof}
At an orthonormal flag all absolute leading Gram determinants are
one. The framed uniformization takes the Hodge directions to the
coordinate vectors and takes the theta generator to one of squared
mean one in the rigidified Gaussian metric on the period torus.
Explicitly its normalization is
\[
  \frac{1}{\operatorname{vol}\mathcal F}\int_{\mathcal F}
  |\theta(z)|^2 e^{-\pi H(z,z)}\,dz=1 .
\]
This is Haar probability on the abelian fibre and therefore
\(\|s_0\|_2=1\) in Proposition~\ref{intjet:arch}. It is part of
the framed uniformization, and changing the scalar of the
generator alone would change this equation. The polarized form
in these coordinates is \(H(z,z)=2x\sum_j|z_j|^2\), so that
\(U=\sqrt{2x}\). The cover metric multiplies the absolute raw
coefficient by \((2\pi x)^{-d_P/2}\) at this flag. In particular
the factor \((\sqrt2 U)^{d_P}\) there becomes
\((2/\pi)^{d_P/2}\), independently of the size of \(x\).
By Lemma~\ref{bfam:minimum} the minimum in the radius in
Proposition~\ref{intjet:arch} satisfies
\(\log^+\lambda^{-1}\le\log R_E+\tfrac12\log(1/\varepsilon)\ll L\).
That proposition, the short coefficient bound above, and
\(\log|\Delta_F|/[F:\mathbb Q]\ll(frL)^2\) give the asserted
estimate at the half maximum. Doubling bounds every frame.

The full cover norm integrates the square of this pointwise
quantity against the probability on the adelic unitary quotient
\(U(W_j)(\mathbb Q)\backslash U(W_j)(\mathbb A)\), writing
the group by restriction of scalars. The component norm uses the
same integrand through the frame map and the component measure
(conditional probability on positive mass). Any such measure of
mass at most one gives the same upper bound by integration.
Embeddings of the other CM orientation are read at their conjugate
good embeddings; the coefficient estimate was simultaneous for both.
This direction of the estimate uses no comparison of two total
adelic volumes.
\end{proof}
\src{src/Formalism/LocalToGlobal/ProjectiveLattices/PetersonNormSubspace.lean}{111}{The squared mean in the framed uniformization};
\src{src/Development/AbelianVarietyTheory/Proofs/DivisorProofs/EvenDomainWeightBound.lean}{68}{The full and component cover norms}.

\paragraph{Descent of the raw rule in characteristic zero.}
One can check that the finite-jet formula gives algebraic values
before identifying it with any classical modular form. Work on
the framed abelian locus over \(F\), with its actual level
structure, eigenframe, theta generator and polarized biextension.

\begin{lemma}\label{bfam:rule}
Assume the column lattice is standard locally at every prime
above the deep level \(M\). The odd-order finite theta sum of Lemma~\ref{newtheta:index}
with the prescribed half-primary division points and harmonic \(P\)
defines a regular rule on these framed points, functorial in
the characteristic-zero base. The complex reading of natural
weight rules on reduced finite-type \(F\)-tests is injective after
extension to \(\mathbb C\), whenever the framed uniformization
surjects onto complex framed points up to the weight action.
\end{lemma}
\begin{proof}
Work locally on the base in logarithm coordinates of the smooth
commutative formal group; they are defined over a \(\mathbb Q\)-algebra
by integrating its invariant differentials, with identity linear term.
On the polarized biextension choose a frame restricting to the
rigidified frames on both axes. Such a frame exists formally
locally: an arbitrary frame at the origin extends over the
completion and its axis restrictions can be divided out.
Use the biextension identification of its pullback by
\(([2],[2])\) with its fourth power. The logarithm of the
discrepancy starts in total degree three. Indeed each first-variable
addition discrepancy vanishes on \(X_1=0,\ X_2=0,\ Y=0\);
the second-variable discrepancy has the analogous property.
Compute dilation by doubling the two variables successively.
Changing the log frame by a degree-\(j\) homogeneous correction
changes the degree-\(j\) dilation discrepancy by \(2^j-4\) times
that correction. Thus all discrepancies of degree at least three
can be removed. Each addition discrepancy of the resulting frame
is now homogeneous of total degree two, by applying simultaneous
dilation to the addition diagram. Its three axis vanishings force
it to be zero.

The Rosati adjunction for an integer \(a\in O\) in this biadditive
frame can differ by an exponential of a bilinear form. In
types \((\tau,\nu)\), let its log coefficient be \(c_a\).
Functoriality for addition and composition of endomorphisms gives
\[
 c_{a+b}=c_a+c_b,\qquad
 c_{ab}=\tau(b)c_a+\nu(\bar a)c_b .
\]
Indeed for composition apply adjunction first to \(a\), with
first input \(bX\), then to \(b\), with second input \(\bar aY\).
If \(\tau=\bar\nu\), this is a derivation on the finite separable
field \(E/\mathbb Q\) with the given coefficient action, hence
zero. If \(\tau\ne\bar\nu\), interchanging \(a,b\) shows
\(c_a\) is a constant multiple of \(\tau(a)-\nu(\bar a)\).
Changing the bilinear log coefficient of the frame removes it.
The remaining ambiguity is Rosati-compatible bilinear. On the
product of \(n\) identical factors it gives the sums of mixed
quadratics killed by harmonic contraction in
Lemma~\ref{intjet:der}. Thus the jet is regular in the base and
independent of these choices. Its evaluation uses only a finite
truncation of these formal constructions and commutes with base
change. The weight action follows by changing the eigenframe in
the polynomial and raising the theta generator's scalar to the
\(n\)-th power.

After an étale base change one may enumerate \(Z_\Lambda\) and
divide each level point by \(N\) in its prescribed half-primary
part. In characteristic zero the torsion of an abelian scheme
is finite étale, and multiplication by \(N\) on the appropriate
primary torsion over geometric points is surjective onto the
prescribed level part. The symmetric normalized lifts on the
odd torsion are algebraic: in the theta group the fibre over
a point is a multiplicative torsor; the two normalizations of
Lemma~\ref{newtheta:index} select a unique lift étale locally.
Indeed first make its \(N^2\)-th power one by taking a root,
then adjust by an \(N^2\)-th root of unity to make inversion
agree with its inverse. The discrepancy is such a root and
squaring is an automorphism on these roots. Uniqueness gives
descent of the lift.

If two choices of division point differ by \(u\), then \(u\)
is \(N\)-torsion supported on the prescribed half above \(M\).
Here the column lattice is standard, so \(u\in K_0\) in the
notation of Lemma~\ref{newtheta:index}, by its period description
at those primes. The Weil pairing between two points in the
same half is trivial: its adjoint involution interchanges the
two relatively prime half ideals, hence their idempotents pair
only across the halves. The translations by \(t\) and \(u\) thus
commute and their normalized product is the lift at \(t+u\).
The vector to which the last translation is applied is fixed by
\(K_0\), by Lemma~\ref{newtheta:index}. This proves independence
of division and of enumeration. The denominator is invertible
in the base, and the level character takes values in \(F\).
The resulting regular scalars agree on the étale double overlap
and descend. On a disconnected base the assertion is checked
componentwise; the same identities hold scheme-theoretically by
the finite étale torsion and the uniqueness of its lifts.

Finally let \(\sum_i c_i f_i\) have vanishing complex reading
at a fixed embedding \(F\to\mathbb C\), with the \(f_i\)
independent over \(F\). On a framed point of any reduced
finite-type test algebra \(R\), it follows that
\(\sum_i c_i\otimes f_i(R)=0\) in \(\mathbb C\otimes_F R\).
Indeed its value is zero at each complex point by naturality,
and the uniformization up to the common weight factor.
This algebra is reduced since
\(F\) is perfect, and its complex points detect zero by the
Nullstellensatz. Expand all \(c_i\) in a basis of their finite
\(F\)-span in \(\mathbb C\). Each coefficient yields an \(F\)-linear
relation among the \(f_i\) on every test, hence is zero. This
proves injectivity.
\end{proof}
\src{src/Formalism/LocalToGlobal/ProjectiveLattices/ThShort.lean}{244}{The faithfully flat evaluation, half-primary division and weight};
\src{src/Formalism/LocalToGlobal/ProjectiveLattices/AlgebraMapCAt.lean}{9}{Reduced finite-type test rings};
\src{src/Development/AbelianVarietyTheory/Proofs/UnitaryExistence/ReadingLinearIndependentAndInjective.lean}{22}{Reading after extension of the coefficient field}.

\paragraph{The operator on row lattices.}
There is a separate numerical reason for the lower bound on the
theta rank. It applies before any denominator or choice of
representatives on the column genus.

\begin{lemma}[Invertibility at the stated theta rank]\label{bfam:gram}
Let \(\mathcal R\) be any set of full lattices of rational rank
\(b=2fr\), self-dual for the fixed rational symplectic polarization.
Set \(m(L,L')=[L:L\cap L']\). Suppose a complex matrix \(C\)
indexed by \(\mathcal R\) has constant diagonal \(A>0\) and
\(|C_{LL'}|\le A\,m(L,L')^{-n/2}\).
If \(n\ge4b+4\), then on bounded lattice-indexed coefficients
\[
     \|A^{-1}C-1\|_\infty\le\pi^2/6-1,\qquad
     \|C^{-1}\|_\infty\le A^{-1}(2-\pi^2/6)^{-1}.
\]
The result holds for coefficients in any complex Banach space
and on any closed subspace preserved by \(C\).
\end{lemma}
\begin{proof}
All self-dual lattices have the same symplectic covolume, hence
\([L':L\cap L']=m(L,L')\). In a rank \(b\) free abelian group
there are \((\ell^b-1)/(\ell-1)\le\ell^b\) subgroups of prime
index \(\ell\), by their kernels in the reduction modulo \(\ell\).
Any index-\(m\) subgroup occurs at the end of a chain with
successive prime indices in a prescribed ordering of the prime
factors of \(m\), by the primary decomposition of the quotient.
There are at most \(m^b\) such chains. The same count holds for
overlattices by ordinary lattice duality. Thus at most \(m^{2b}\)
choices of \(L'\) have this index, by first choosing \(L\cap L'\)
and then its overlattice. Index one gives only \(L\). The sum of
absolute off-diagonal entries in each normalized row is bounded
by \(\sum_{m\ge2}m^{2b-n/2}\le\sum_{m\ge2}m^{-2}\).
The norm-convergent Neumann series proves both conclusions,
including on the indicated subspaces.
\end{proof}

Here is the finite vacuum calculation giving the entry estimate.
For unshifted row lattices \(L,L'\), put
\(I=L\cap L'\). For any locally free rank-\(n\) column lattice
\(\Lambda\) self-dual for its Hermitian form, write
\(A_L=\Lambda\otimes_O L\) for the corresponding tensor lattice
with its trace symplectic form. Then
\[
 A_L\cap A_{L'}=\Lambda\otimes_O I,\qquad
 [A_L:A_L\cap A_{L'}]=m(L,L')^n .
\]
The first identity is flatness and the second follows on choosing
local bases at the finite primes. The two tensor lattices are
self-dual for the symplectic form since the column pairing
is perfect over \(O\). In the Heisenberg representation take
unit vacuum vectors \(v,v'\) with the prescribed characters on
these two lattices. Distinct translates of \(v'\) by representatives
of \(A_L/(A_L\cap A_{L'})\) are orthogonal: on \(A_{L'}\) their
eigencharacters differ by the symplectic commutator, whose
annihilator is precisely \(A_{L'}\). Their inner products with
\(v\) all have the same absolute value, because \(v\) is an
eigenvector on \(A_L\). Bessel's inequality gives
\(|\langle v,v'\rangle|\le m(L,L')^{-n/2}\).
This argument works with any translates of the two vacua, since
translation just changes their eigencharacters.

For the actual finite marks write \(\mathfrak m\) for the nonzero
proper odd ideal prime to \(\bar{\mathfrak m}\), and \(L_*\) for
the reference row lattice. The row set consists of self-dual \(L\)
equal to \(L_*\) at primes dividing \(\mathfrak m\bar{\mathfrak m}\).
The column lattice is standard there. Set
\(B_L=\Lambda\otimes_O(\mathfrak m\bar{\mathfrak m}^{-1}L)\).
The twist ideal times its conjugate is \(O\); hence these \(B_L\)
are again self-dual and their relative index remains \(m(L,L')^n\).
Choose a basis of \(L_*/\mathfrak m L_*\). The marked vector
\(\phi_L\) is the sum of translates of the unit vacuum on \(B_L\).
They are indexed by \(k\in J\subset{\rm GL}_r(O/\mathfrak m)\)
and weighted by the unitary level character. In fact the quotient
used to choose the translating tuples is the reduction modulo
\(\mathfrak m\) of \(\Lambda\otimes_O(2\bar{\mathfrak m}L)\).
Since \(2\bar{\mathfrak m}\) is locally a unit on \(\mathfrak m\),
this gives precisely the \(n\) marked rows modulo \(\mathfrak m L_*\).
The first \(r\) rows give the basis transformed by \(k\), and
the remaining rows vanish modulo \(\mathfrak m\). There is one
class for each such \(k\).

Use the common vacuum semicharacter on intersections. For
translates by \(y,y'\), nonzero inner product implies triviality
of the commutator of \(y-y'\) on \(B_L\cap B_{L'}\), hence
\(y-y'\in (B_L\cap B_{L'})^\#=B_L+B_{L'}\).
Here the duality identity follows by dualizing the sum of the
two self-dual lattices. At every \(\mathfrak p\mid\mathfrak m\)
this sum has rows in \(\mathfrak m L_{*,\mathfrak p}\).
The marked residues coincide, so only pairs with the same \(k\)
contribute. Their diagonal squared norm is \(|J|\); the mixed
inner product is at most \(|J|m(L,L')^{-n/2}\).
This proves precisely the hypotheses of Lemma~\ref{bfam:gram}
for the positive weighted class sum, with \(A=|J|\) times the
total weight. No lower bound on that mass is needed for invertibility.

The same calculation detects the arithmetic group of the marks.
For \(\gamma\in U(h_0)(E^+)\) replace \(B_{L'}\) by \(\gamma B_{L'}\);
the vacuum semicharacter is preserved because it is defined by the
polarization, the fixed \(O\)-scalar \(w_0\) with \(w_0+\bar w_0=1\),
and the unitary action. Nonzero overlap now gives
\(y-b=\gamma(y'-b')\) for some \(b\in B_L,\ b'\in B_{L'}\),
using the same duality identity. At \(\mathfrak p\mid\mathfrak m\)
both adjusted tuples have first \(r\) rows forming an integral
basis of \(L_{*,\mathfrak p}\), by reduction modulo the maximal
ideal. Thus \(\gamma\) preserves the localized reference lattice
and its residue maps the basis marked by \(k'\) to that marked by
\(k\). Its residue is \(k(k')^{-1}\in J\). This proves vanishing
on the graph orbits outside the subgroup specified by these
congruences.

One also obtains the finite vanishing along a nonzero isotropic
row line \(E\xi\). For completeness take a prime
\(\mathfrak p\mid\mathfrak m\), put \(a={\rm ord}_{\mathfrak p}
(\mathfrak m)>0\), and work at this prime and its distinct conjugate.
The trace symplectic pairing identifies \(L_{*,\bar{\mathfrak p}}\)
with the additive character dual lattice of \(L_{*,\mathfrak p}\)
in \(E_{\bar{\mathfrak p}}^r\). This follows by testing the
global self-duality locally on the two conjugate idempotents.
Scale \(\xi_{\bar{\mathfrak p}}\) to a primitive vector \(u\)
of \(L_{*,\bar{\mathfrak p}}\). The element
\(\varpi_{\bar{\mathfrak p}}^{-a}u\) belongs to the shifted row
lattice there but lies outside the additive dual of \(L_{*,\mathfrak p}\).
For any marked tuple \(y\) its first \(r\) rows form a basis at
\(\mathfrak p\). Some integral multiple (in the local integer
ring) of \(\varpi_{\bar{\mathfrak p}}^{-a}u\) therefore pairs
nontrivially under the additive character with one of those rows.
Put this vector in the corresponding column position of a tuple
\(\delta\in B_L\otimes\widehat{\mathbb Z}\) on \(E^n\otimes E\xi\),
and zero at the other components. Equivalently one can choose a
rational \(\delta\) giving the same character: the intersection of
a full lattice with a rational subspace spans a lattice in that
subspace, and its completion is the intersection of the completed
lattice with the completed subspace, by a basis adapted to the
saturated intersection. Approximation modulo a finite-index
sublattice then preserves the character.
On this subspace both the alternating form and the vacuum quadratic
phase vanish, since \(h_0(\xi,\xi)=0\). The Heisenberg translations
on it commute and their action on the marked translate has the
character of this pairing. For each of the finitely many labels
choose such a \(\delta\) with eigenvalue \(\eta\ne1\). The product
of the commuting operators \(T_\delta-\eta\) kills the sum of
marked translates and has nonzero augmentation
\(\prod(1-\eta)\). Hence its image in the coinvariants of the
isotropic row subgroup vanishes. This supplies the vanishing for
non-graph terms whose evaluation factors through those coinvariants.
\src{src/Development/AbelianVarietyTheory/Proofs/SiegelThetaTransformation/ThetaTestFunctionNormAndOverlapBound.lean}{133}{The localized congruence for a surviving pair of labels};
\src{src/Development/AbelianVarietyTheory/Proofs/SiegelThetaTransformation/HermitianThetaLevelDataFormsSystem.lean}{60}{The off-diagonal sum at rank \(8fr+4\)};
\src{src/Development/AbelianVarietyTheory/Proofs/SiegelThetaTransformation/ThetaTestFunctionNoCoinvariantsIsotropic.lean}{216}{Isotropic row coinvariants of the same marked tests}.

\paragraph{Which definite Siegel--Weil formula is needed.}
There is a way to use the ordinary adelic theorem without comparing
the volumes of different genera of self-dual Hermitian lattices.
Use just the genus of \(\Lambda_*=O^n\) in the positive definite
space \(V=E^n\), and require equality with the standard lattice
at a prescribed integer \(B_0\) divisible by the deep level and
the excluded rational primes. Every lattice used in this smaller
family is among the permitted columns. Its positive class mass
still gives the strict inequality in Lemma~\ref{bfam:gram}.

\begin{lemma}[Adelic realization of a lattice family]\label{bfam:genus}
Write \(H=U(V)\) over \(E^+\) and suppose \(n>4r\).
For each such column lattice let
\(\varphi_\Lambda\) be a function on \(V(E^+)^r\) (thus on
\(r\) columns in \(E^n\)). Assume the following conditions for
some positive integer \(R\).
\begin{enumerate}
\item The function is supported in \(R^{-1}\Lambda^r\) and
invariant under \(R\Lambda^r\).
\item If two lattices coincide at all primes over \(RB_0\),
their functions agree on the intersection of their indicated
supports \(R^{-1}\Lambda^r\cap R^{-1}\Lambda'^r\).
\item For a rational isometry \(\gamma\) integral and congruent
to the identity modulo \(B_0\) at its prime divisors,
\(\varphi_{\gamma\Lambda}(\gamma x)=\varphi_\Lambda(x)\).
\end{enumerate}
Then the Schwartz functions obtained by completion form the
translates of one finite adelic Schwartz function in this genus.
For a fixed \(H(E^+\otimes\mathbb R)\)-invariant archimedean
Schwartz function and the splitting in which \(H\) acts by
linear substitution in the doubled Schrödinger model, the
genus average of the theta functions with stabilizer weights,
divided by the sum of those weights, is the anisotropic
Siegel--Weil Eisenstein value for the standard lattice test.
The same is true of the genus average of the corresponding
Eisenstein values themselves.
\end{lemma}
\begin{proof}
Let \(\mathcal U\subset H(\mathbb A_f)\) be unrestricted away
from \(B_0\), and the principal congruence subgroup of the
standard lattice at its divisors. For \(g\in\mathcal U\)
consider \(\Lambda'=g\widehat\Lambda\cap E^n\). Write
\(\widehat\varphi\) for the function on the finite adeles
defined by the same finite quotient
\(R^{-1}\Lambda^r/R\Lambda^r\), extended by zero.
We claim
\(\widehat\varphi_{\Lambda'}(g x)=\widehat\varphi_\Lambda(x)\).
It suffices to check this on \(R^{-1}\widehat\Lambda^r\).
Choose a rational unitary \(\gamma\) sufficiently close to \(g\)
at all primes over \(RB_0\). Weak approximation here follows
from the Cayley transform between the skew-adjoint Lie algebra
and a dense open in \(H\); the latter is locally dense at every
finite place, since its complement is a proper algebraic closed
set in a smooth connected group. Linear weak approximation in
the Lie algebra proves the assertion. Thus \(\gamma\Lambda\)
and \(\Lambda'\) agree at those primes, \(\gamma\) meets the
congruences of the third hypothesis, and \(\gamma x,g x\)
have the same residues modulo the local period lattice there.
The lattice \(R^{-1}(\gamma\Lambda\cap\Lambda')^r\) maps onto
these residues by the Chinese remainder theorem. Choose a
rational \(z\) in it representing them. Outside \(R\) the
support lattice and period lattice coincide. Consequently the
second hypothesis compares \(\varphi_{\Lambda'}(z)\) and
\(\varphi_{\gamma\Lambda}(z)\), and the third identifies the latter
with \(\widehat\varphi_\Lambda(x)\). This proves the claim.

The same weak approximation at \(B_0\) gives
\(H(\mathbb A_f)=H(E^+)\mathcal U\).
Put \(K=\operatorname{Stab}_{\mathcal U}(\widehat\Lambda_*)\).
The compact right cosets on the finite double quotient
\(H(E^+)\backslash H(\mathbb A_f)/K\) correspond to the
lattice orbits with the indicated congruences. Representatives
can all be taken in \(\mathcal U\). Right invariance by \(K\)
follows from the claim. On each coset the compact archimedean
group integrates the invariant archimedean factor, and the
remaining mass is proportional to the reciprocal order of
the stabilizer. Indeed its chart is a quotient of
\(H(E^+\otimes\mathbb R)\times K\) by this finite group;
the numerator Haar volume is common to all cosets. Dividing by
the total mass gives the theta integral for Haar probability.
The doubled split group has Witt index \(r\); the source
dimension is \(n>4r\). In the notation of Gan--Qiu--Takeda the
anisotropic formula is at \(s=(n-r)/2>0\), with unitary
constant one for the probability theta integral
\cite[\S3.1 and Theorem~7.1]{ganQiuTakeda}. Finally each
finite transport has the same Siegel section: the \(H\)-action
commutes with the doubled group and linear substitution fixes
evaluation at zero. This proves the last assertion as well.
\end{proof}
For an archimedean polynomial insertion one takes the invariant
sum over a unitary orthonormal basis of the column representation
in Proposition~\ref{bfam:harm}, paired with its conjugate on
the second copy. This is the identity endomorphism of that
representation, so the resulting doubled Schwartz function is
invariant under the compact column group. Lemma~\ref{bfam:genus}
applies just as it does to the Gaussian. We verify the mixed
finite equivariance directly.
\src{src/Development/AbelianVarietyTheory/Definitions/SpanningFamilies/FirstCoordinateAdmissible.lean}{139}{Equivariance, periodicity and finite locality of the column family};
\src{src/Development/AbelianVarietyTheory/Proofs/SiegelThetaTransformation/AdmissibleDoubledHermitianThetaFamily.lean}{413}{Mixed vacuum comparison for the marked doubled test}.

\begin{lemma}[The finite phase for the projected marked vacua]\label{bfam:phase}
Use the finite adelic Schrödinger representation of the trace
symplectic space underlying \(V\otimes_E E^r\). For each row
lattice \(L_i\), \(i=1,2\), start with any unit vacuum \(v_i\) for
the standard column and the shifted row lattice
\(D_i=\mathfrak m\bar{\mathfrak m}^{-1}L_i\). For the column
\(\Lambda\) take the unit normalization \(u_i^\Lambda\) of the
orthogonal projection of \(v_i\) to the vacuum of
\(A_i^\Lambda=\widehat\Lambda\otimes_{\widehat O}\widehat D_i\).
The vacuum character at each lattice is the prescribed
semicharacter defined by \(w_0\). Form the marked sums
\(\phi_i^\Lambda\) above from these \(u_i^\Lambda\). For
\(h\in H(\mathbb A_f)\) preserving the marked residues (in
particular for \(h\in\mathcal U\) of the preceding proof),
\[
 W_{\phi_1^{h\Lambda},\phi_2^{h\Lambda}}(hx)
       =W_{\phi_1^\Lambda,\phi_2^\Lambda}(x),\qquad
 W_{u,v}(x)=\langle \rho(x)u,v\rangle .
\]
Here \(h\Lambda\) means the lattice with completion
\(h\widehat\Lambda\), and \(\rho\) denotes the Weyl operators.
In particular the finite doubled tests used in the ordinary
adelic theta identity are the translates of a single test for
each pair of row lattices, with their original complex coefficients.
\end{lemma}
\begin{proof}
Choose complementary rational Lagrangians in the underlying
symplectic space. Use dual Haar measures for the finite additive
character \(\psi_f\) and the pairing of the two Lagrangians. With
the convention \(\rho(a,b)u(t)=\psi_f(b(t+a/2))u(t+a)\) and the
inner product linear in its first argument, the map to the doubled
diagonal model is
\[
 W_{u,v}(a,b)=\int u(t+a/2)\overline{v(t-a/2)}\,\psi_f(bt)\,dt .
\]
Indeed one first changes variables from the two position coordinates
to their difference and midpoint, then Fourier transforms the
midpoint. The Jacobian of the first change is of absolute value
one, also at dyadic places. A change of basis for the dual
coordinate changes the corresponding dual Haar. Fourier inversion
and this change of variables give
\[
 \int W_{u,v}\overline{W_{s,t}}=
             \langle u,s\rangle\overline{\langle v,t\rangle}.
\]
The integrals here are integrals of Schwartz functions. In
particular \(W_{u,v}(0)=\langle u,v\rangle\).
In the paper's rational doubled coordinates the symbol is
\(a\,W_{u,v}(-x)\), where \(a\ne0\) is the common change of Haar
normalization. In fact translation of the two inputs by
\(y_1,y_2\) translates the position argument by \(-(y_1-y_2)\)
in those coordinates. Set \(y_1=-x,\ y_2=0\) and evaluate at
the origin. The remaining half-pairing has zero rational trace:
it is a purely imaginary scalar times the real Hermitian norm
of \(x\). Thus the value of the doubled symbol at \(x\) equals
its origin value on \(\rho(-x)u,v\). The origin pairing is the
displayed inner product by partial Fourier transformation. The
reflection commutes with every linear column isometry.

Write \(\epsilon(a)\rho(a)\) for the action fixing a vacuum on a
self-dual lattice \(A\). The formula for \(\epsilon\) depends on
the vector and \(w_0\), so its restrictions to intersections
agree. Explicitly it is the half quadratic character of the trace
pairing of \(a\) with \(w_0 a\), with the sign determined by
\(\psi_f\). On \(A\) it cancels the Weyl multiplication cocycle:
\(w_0+\bar w_0=1\) and the pairing of \(A\) with \(w_0A\) is
integral. The integral of \(\epsilon(a)\rho(a)\) against Haar
probability on \(A\) is the orthogonal projection \(P_A\) to its
fixed space. For a vacuum \(v\) on \(A_*\), \(W_{v,v}\) vanishes
outside \(A_*\): conjugation by \(b\in A_*\) multiplies the
coefficient by the commutator character and \(A_*^\#=A_*\).
Thus
\[
 \|P_Av\|^2=\langle P_Av,v\rangle
      =[A:A\cap A_*]^{-1}.
\]
The norm is positive and the unit projection has positive real
inner product with \(v\). This integral is also the normalized
finite average of the distinct translates used on the rational
lattice: the orbit on a Schwartz vector is finite and each fibre
of its orbit map has equal measure.

For vacua \(u_i\) on \(A_i\), the same commutator, tested on
\(A_1\cap A_2\), gives support in \(A_1+A_2\). At
\(x=a_1+a_2\), \(a_i\in A_i\), move \(\rho(a_1)\) to the
first vacuum and \(\rho(a_2)\) to the second. The result is
\(W_{u_1,u_2}(0)\) times a unit scalar determined by
\(\epsilon(a_i)\) and their Weyl cocycle. In particular its
absolute value is constant throughout the support. The preceding
Fourier identity on the diagonal gives
\[
 |W_{u_1,u_2}(0)|={\rm vol}(A_1+A_2)^{-1/2}.
\]
This volume is independent of the self-dual column: each
self-dual \(A_i\) has volume one and the relative index here
is \([D_2:D_1\cap D_2]^n\).

To compare the \emph{arguments} for columns \(\Lambda,\Lambda_*\),
observe that the intersection of the supports is
\[
 (\widehat\Lambda\cap\widehat\Lambda_*)\otimes_{\widehat O}
                        (\widehat D_1+\widehat D_2).
\]
This can be checked in a local basis of the row sum. Every
element of this intersection therefore has a decomposition
\(a_1+a_2\) with the two \(a_i\) in the corresponding vacuum
lattices for both columns. Its scalar relative to the value at
zero is identical for the two symbols. Their Fourier inner
product is consequently the positive volume of this intersection
times the product of the two values at zero, with the second
conjugated. On the other hand the Fourier identity identifies
this inner product with
\(\langle u_1^\Lambda,v_1\rangle
 \overline{\langle u_2^\Lambda,v_2\rangle}>0\).
The two values at zero have equal modulus, so they are equal.
This applies separately to every column and does not impose a
choice of phase on either initial \(v_i\).

An isometry \(h\) of the column carries a decomposition
\(a_1+a_2\) to \(ha_1+ha_2\). It preserves the half quadratic
character and the Weyl cocycle. Equality at zero now gives the
assertion for the unmarked vacua. For a marked summand translate
by \(y_i\in\widehat\Lambda\otimes 2\bar{\mathfrak m}\widehat L_i\).
Its doubled symbol is the unmarked symbol at a translate of \(x\)
multiplied by the Weyl phase of \(x,y_1,y_2\); the latter is
also invariant under \(h\). The lifts \(hy_i\) give the prescribed
residues for the new column. Any two lifts of the same residue
differ by an element of \(2A_i^{h\Lambda}\), as may be checked at
the primes of \(\mathfrak m,\bar{\mathfrak m}\) and their
complement. Its vacuum character is trivial, and its half
pairing with the chosen lift is integral since both lifts lie in
\(2(h\widehat\Lambda\otimes\widehat L_i)\).
Thus each marked translate is independent of this choice.
Rational lifts with these residues exist by the Chinese remainder
theorem. Summing with the given weights on \(J\times J\) proves
the formula.
\end{proof}
The proof uses the finite additive character trivial on
\(\widehat{\mathbb Z}\) obtained from the rational trace pairing.
It gives the finite diagram also at \(2\), with the specified
semicharacter. In particular the normalized class average of
these genuine doubled symbols is exactly the finite input of
Lemma~\ref{bfam:genus}. The choices of phases for the initial
row vacua are allowed to depend on \(L_i\).
\src{src/Development/AbelianVarietyTheory/Definitions/SpanningFamilies/FirstCoordinateAdmissible.lean}{81}{The sign in the doubled position coordinate};
\src{src/Development/AbelianVarietyTheory/Proofs/SiegelThetaTransformation/DoubledHermitianTestFunctionZeroAndNorm.lean}{275}{The common origin normalization};
\src{src/Development/AbelianVarietyTheory/Definitions/SpanningFamilies/IsSystem.lean}{148}{Vacuum semicharacter and the normalized projection};
\src{src/Development/AbelianVarietyTheory/Proofs/SiegelThetaTransformation/DoubledVacuumEqHermitianCharacter.lean}{554}{Constancy of the mixed scalar across columns}.

In the doubled model a rational maximal isotropic subspace of
\(W\oplus(-W)\) either is the graph of an isometry, or contains an
isotropic row line in \(W\oplus0\). Indeed if its first projection
has a kernel, its second projection is orthogonal to that kernel,
and hence has a nonzero kernel as well. If both projections are
invertible, isotropy is precisely the isometry equation. On a
graph the finite evaluation at zero is the inner product with the
second vacuum transported by the row isometry, since the base
graph has functional \(W_{u,v}(0)\). For a nongraph the evaluation
is invariant under Weyl translations in the indicated row line
tensored with \(V\). In the polarization of this maximal isotropic
subspace they are multiplication by additive characters of the
position variable and each such character is one at zero. Thus
the augmentation argument above kills this orbit term, including
after archimedean polynomial insertion. This accounts for all
rational isotropic subspaces.

\begin{corollary}[The Fourier mark on the identity cover]\label{bfam:mark}
At the deep level used for the raw spanning calculation the
auxiliary map of that calculation is the identity. Its auxiliary
ideal is \(O\), its integer is one and its auxiliary character is
trivial. Thus \(\mathfrak m\) is the same half of the level on
both sides. The subgroup \(J\) of residues is the identity:
its entries differ from the identity by reductions of elements
of \(\mathfrak m\) modulo \(\mathfrak m\). The algebraic
Fourier mark on the half-primary torsion selects exactly the
single translated shifted vacuum with the first \(r\) basis
rows prescribed and the remaining rows zero.
\end{corollary}
\begin{proof}
Put \(A=\widehat\Lambda\otimes\widehat L\) and
\(D=\widehat\Lambda\otimes\mathfrak m\bar{\mathfrak m}^{-1}
\widehat L\), and use a subscript \(*\) for the standard column.
The alternating commutator identifies \(A/(A\cap D)\) with
the character dual of \(D/(A\cap D)\): each induces a character
on the other group, and self-duality of both lattices gives
trivial kernel on both sides. The second group is supported on
\(\bar{\mathfrak m}\), where it is
\((\bar{\mathfrak m}^{-1} L/L)^n\); at those primes the column is
standard. It is the group \(Q\) of the \(n\) half-level torsion
points. To see the identification, use the level isomorphism
on \(M\)-torsion and the two Chinese remainder idempotents of
\(O/(M)\), since \((M)=\mathfrak m\bar{\mathfrak m}\).
The dual group similarly has residues in \((L/\mathfrak m L)^n\).
In these identifications the Fourier coefficient of the chosen
algebraic mark is precisely the pairing character for those
residues; this is the defining property of the dual marking.
One can lift its tuple by \(y\in\widehat\Lambda\otimes
2\bar{\mathfrak m}\widehat L\), since \(2\bar{\mathfrak m}\)
is a unit at \(\mathfrak m\). Let \(v\) be a unit unshifted
vacuum for the standard column and put \(z=P_Av\). The restriction
of \(\epsilon(d)\rho(d)\) to the intersection fixes \(z\).
Consequently averaging these operators on \(z\) amounts to averaging
over \(Q\). Conjugation by \(\rho(y)\) multiplies the summands
by their pairing character. Moreover \(\rho(y)z=z\), because
\(y\in2A\). With the character taken from this conjugation, the
normalized Fourier average is therefore \(\rho(y)P_D P_Av\).
Taking the opposite convention for the additive character conjugates
both the definition of the dual mark and the commutator.
The odd-order canonical Heisenberg lifts agree on \(Q\) with this
action on \(z\). Indeed the lifts on its isotropic subgroup are
additive, inversion acts by negation, and a change of such a lift
is a character of order dividing two on the odd group.

The normalization here can be compared on \emph{all} rows at
once. Pick the initial unit shifted vacuum to be the normalized
\(P_{D_*}v\). At the primes over \(M\), \(A=A_*,D=D_*\);
elsewhere \(D=A,D_*=A_*\). The local projectors at distinct
places commute, so on \(v\) we have
\(P_DP_Av=P_DP_{D_*}v\).
The product statement follows by factoring a vacuum for \(A_*\)
into local vacua. For completeness the vacuum line of a self-dual
lattice is one dimensional: the calculation of \(W_{u,v}\) above
gives \(|\langle u,v\rangle|=1\) for any two unit vectors fixed
with its character. The restricted tensor of local vacua supplies
one such vector, hence any choice differs only by a unit scalar.
Projector norms in Lemma~\ref{bfam:phase} now show that the
coefficient of the unit shifted marked vacuum is the positive
number
\[
 [D_*:D_*\cap A_*]^{-1/2}\,
 [D:D\cap D_*]^{-1/2}.
\]
The first index is \(|O/\mathfrak m|^{nr}\). The second is
\([\widehat\Lambda:\widehat\Lambda\cap\widehat\Lambda_*]^r\),
by using local bases of the shifted row lattice. Both are
independent of the row and its component. Changing the original
shifted unit vacuum for a fixed row multiplies all its projected
unit vacua by the same phase. Finally, the residue conditions on
the last \(n-r\) rows on both sides of the comparison are zero
modulo the identical ideal \(\mathfrak m\).
\end{proof}
\src{src/Formalism/LocalToGlobal/ProjectiveLattices/ThShort.lean}{300}{Identity auxiliary map on the deep model};
\src{src/Development/AbelianVarietyTheory/Definitions/SpanningFamilies/IsBaseClassLevelAdapted.lean}{78}{The two half-level ideals and the dual residue marking};
\src{src/Development/AbelianVarietyTheory/Proofs/SiegelThetaTransformation/LeSupOfCoprime.lean}{100}{Composition of the unshifted and shifted projections}.

\paragraph{Integral change by a flat character.}
Fix an algebraic component \(Z\) of the cover over \(F\).
For a rational trivialization \(t\) of a flat automorphic line on
\(Z\), suppose its pointwise metric is constant \(r_\sigma>0\)
on each complex fibre, with conjugate embeddings giving the same
constant. Suppose \(\rho_w\) is a lower bound for its vertical
orders at every test over the finite prime \(w\) of \(F\), in the
normalization which gives the order \({\rm ord}_w(a)\) to
\(a\in F^\times\). Assume \(\rho\) has finite support and is
integral off a finite set \(B\) of rational primes. Put \(h=[F:\mathbb Q]\)
and
\[
 \delta(t)=h^{-1}\left(\sum_w\rho_w\log Nw-
                         \sum_{\sigma:F\hookrightarrow\mathbb C}
                               \log r_\sigma\right).
\]
There is \(a\in F^\times\) such that \(at\) is integral at all
the tests and
\begin{equation}\label{bfam:flat}
 \max_\sigma\log\|at\|_\sigma
 \le-\delta(t)+\frac{\log|\Delta_F|}{2h}+\sum_{\ell\in B}\log\ell .
\end{equation}
Indeed take \(\mathfrak b=\prod_w w^{\lceil-\rho_w\rceil}\).
Every element of this ideal gives integrality. The rounding cost in
\(\log N\mathfrak b\) over \(\ell\in B\) is at most
\(\sum_{w\mid\ell}\log Nw\le h\log\ell\).
In the Minkowski embedding use the intervals and complex discs
with radii \(R/r_\sigma\) at real and complex places. Their volume
is \(2^{r_1}\pi^{r_2} R^h/\prod_\sigma r_\sigma\).
The covolume of \(\mathfrak b\) is
\(2^{-r_2}\sqrt{|\Delta_F|}\,N\mathfrak b\).
Taking \(R^h=\sqrt{|\Delta_F|}\,N\mathfrak b\prod_\sigma r_\sigma\)
achieves Minkowski's bound, including the equality case by compactness.
It gives \eqref{bfam:flat}.

If \(t\) is a nowhere vanishing section on the characteristic-zero
component, multiplication is a field-linear isomorphism to the
twisted cusp space there whenever the flat trivialization and
its inverse preserve cusp boundary conditions. Extension by zero
from the component and restriction back are inverse on its supported
sections. Thus any field-spanning family of global sections restricts
to a spanning family on \(Z\) under this component decomposition.
Multiplication by \(at\) preserves the vertical inequalities and
multiplies the component Petersson norm by \(|\sigma a|r_\sigma\):
this follows under the integral defining that norm, since the
absolute metric of the multiplier is constant on the component.
The inverse of \(at\) is needed only over \(F\). In particular no
division by a component mass or by a generator of a fractional
ideal is needed in the integral lattice.
\src{src/Development/AbelianVarietyTheory/Proofs/DivisorProofs/ExistsShortIntegralFlatLineTwistingSection.lean}{203}{Minkowski rescaling for fractional vertical orders};
\src{src/Development/AbelianVarietyTheory/Proofs/DivisorProofs/ExistsSpanningFamilyMulTwistingSection.lean}{97}{Component projection and multiplication on the field-spanning family}.

\paragraph{Passage from the marked family to the cover.}
Fix a good embedding of \(F\) and a component
\(k\). Let \(L_k\) be its row lattice after transport of its level
marking into the common rational Hermitian space. Write
\(\Theta_{\Lambda,P}(L_k,x)\) for the corresponding analytic
theta jet on a frame \(x\), using the marked vector of
Corollary~\ref{bfam:mark}, and write \(f_{\Lambda,N,P}\) for the
algebraic rule. Identify the polynomial variables by the
eigenframe and the framed periods. The required reading has the form
\begin{equation}\label{bfam:reading}
  f_{\Lambda,N,P}(k,x)=C_k\,u_{\Lambda,N}\,
                   \Theta_{\Lambda,P}(L_k,x).
\end{equation}
Here \(C_k\ne0\) is independent of \(\Lambda,N,P,x\); the
nonzero \(u_{\Lambda,N}\) is independent of \(k,P,x\).
A different normalization for the initial vacuum for \(L_k\)
can be absorbed in \(C_k\).

\begin{proposition}[The theta reading on a polarized period chart]
\label{bfam:unitread}
Use the principal PEL period chart of Corollary~\ref{bfam:metric}
on one connected component. Use its rigidified identification
\(\mathcal L=(1,w_0)^*\mathcal B\) with the pullback of its
polarized biextension, its given half-level marking and its dual,
and its labelled eigenframe as the tangent coordinates. Use the
identity auxiliary map of Corollary~\ref{bfam:mark}. For the raw
section of Lemma~\ref{newtheta:index}, \eqref{bfam:reading} holds
on unit flags up to a common factor of constant positive absolute
value in place of \(C_k\). With the shifted standard vacua chosen
as in Corollary~\ref{bfam:mark}, one may take
\[
        u_{\Lambda,N}=|Q|^{-1/2}
           [\Lambda:\Lambda\cap O^n]^{r/2}.
\]
For the convention that divides the raw rule by \(\beta_H\),
multiply this number by \(\sigma(\beta_H)^{-1}\).
On the good framed chart, \eqref{bfam:reading} holds on the
connected component with a single \(C_k\ne0\).
\end{proposition}
\begin{proof}
Write \(\Pi\) for the row period lattice and \(g=fr\).
In the Appell--Humbert frame over the Lie algebra the automorphy
factor of \(\mathcal L\) at a period \(v\) has unitary
semicharacter \(\chi(v)\) and Hermitian form \(H\).
The pullback by \((1,w_0)\) of the addition biextension has
Hermitian form \(H(z,w_0z')+H(w_0z,z')=H(z,z')\), by Rosati
adjunction. Its semicharacter is the quotient of the characters
at \((1+w_0)v,v,w_0v\), namely
\(\exp(\pi i E_L(v,w_0v))\) by the semicharacter law. Uniqueness
of the Appell--Humbert data of a rigidified line gives this same
semicharacter for \(\mathcal L\). Its values on periods are signs.

Choose a symplectic \(\mathbb Z\)-basis of \(\Pi\) and denote the
two period matrices in the unit tangent frame by \(P_0,Q_0\).
Set \(h=P_0^{-1}\), \(\Omega=h Q_0\), \(Y=\Im\Omega\). The
Riemann bilinear relations give \(Y>0\), \(\Omega^t=\Omega\)
and \(H(z,z)=(hz)^tY^{-1}\overline{hz}\). The signs in this
basis specify a half characteristic \((\alpha,\eta)\). Put
\[
 T(z)=e^{\pi(hz)^tY^{-1}(hz)/2}
          \sum_{\ell\in\mathbb Z^g+\alpha}
           e^{\pi i\ell^t\Omega\ell+2\pi i\ell^t(hz+\eta)}.
\]
Its multipliers follow by shifting the summation index on imaginary
periods and multiplying each term on real periods. The signs
produce the specified characteristic. The same period laws
force every holomorphic section to be a scalar times \(T\):
remove the first exponential, take the Fourier series in the real
periods, then use the other periods to recover every coefficient
from one. This also proves \(T\ne0\).

For a unit finite vacuum for \(A_*\), the \(n\)-fold Schrödinger
theta in these coordinates is
\[
       d_L(\det h)^{n/2}\prod_{j=1}^n T(z_j)
\]
with a fixed \(d_L\ne0\). In the basis adapted to \(\Pi\) the
integral translations and their dual characters force the finite
test to be supported on the translate
\(\alpha+\widehat{\mathbb Z}^{g}\) in each row; its values there
are fixed multiples of the character prescribed by \(\eta\).
Restriction of this function to the rational points gives exactly
the displayed characteristic series. The determinant is the
weight-\(n/2\) factor for the change \(z\mapsto hz\).
To compare to a fixed initial rational symplectic basis, it suffices
to change basis by a linear substitution, a symmetric shear, or
Fourier exchange. These generate the rational symplectic group by
symplectic elimination. Linear substitution reindexes the rational
sum. For a shear the archimedean and finite quadratic characters
at every rational point are inverse by the product formula. For
exchange, apply Poisson summation with dual Haar measures to the
finite test times the archimedean Gaussian with linear argument.
Completing the square gives in each row the factor
\(\det(-i\Omega)^{-1/2}
e^{-\pi i\zeta^t\Omega^{-1}\zeta}\) for the Fourier transform
with shifted argument \(\zeta\). The quadratic exponential is
absorbed by the change of the holomorphic Gaussian prefactor in
\(T\), and the half determinant by the change of \(\det h\).
Thus the new finite vacuum is multiplied by a fixed scalar
independent of the period point; choices of square root give
only a constant. This is the period-coordinate version of the
Schrödinger theta transformation.

The squared mean of \(T\) on the polarized torus is
\(\det(2Y)^{-1/2}\). In fact write \(hz=t+\Omega s\) in
the unit period cube. The integral over \(t\) kills all unequal
Fourier indices. The remaining integrand, including the metric
\(e^{-\pi H(z,z)}\), gives the translates of
\(e^{-2\pi s^t Y s}\); integrating the sum over the \(s\)-cube
gives its integral over \(\mathbb R^g\).
The squared mean of the \(n\)-fold Schrödinger product is
therefore
\[
    |d_L|^2|\det h|^n\det(2Y)^{-n/2}
      = |d_L|^2\bigl(\det H/2^g\bigr)^{n/2}.
\]
At unit flags \(\det H=(2x_0')^g\), where \(x_0'>0\) denotes
the common absolute polarization scalar called \(x\) in
Corollary~\ref{bfam:metric}; equivalently one may choose
orthonormal flags first and extend by the flag character.
The determinant is constant. The geometric generator has squared
mean one. Hence its \(n\)-fold product is the Schrödinger product
times a scalar of constant positive modulus, independent of all
column data and polynomial insertions.

It remains to check what the geometric translations and
differentiations do to this one product. In these coordinates
put
\[
  (U_t^{(m)}s)(z)=e^{-\pi mH(z+t/2,t)}s(z+t).
\]
It sends \([N]^*s=s(Nz)\) under \(U_t^{(N^2)}\) to
\((U_{Nt}^{(1)}s)(Nz)\). This pullback equality in the
Appell--Humbert frames follows since the factors of automorphy
agree and the isomorphism is rigidified at zero. When
\(N^2t\in2\Pi\), inversion sends this translate to its inverse
and its \(N^2\)-th power acts trivially: the residual period
sign at \(N^2t\) is one. Thus it is exactly the normalized
geometric lift of the rule. Such representatives can always
be chosen since \(N\) is odd. On an individual summand in
period coordinates \(u=a+\Omega b\) the finite Heisenberg
operator on a rational coefficient function is
\[
   v(s)\longmapsto
             e^{2\pi i a^t(s-b/2)}v(s-b).
\]
Reindex \(s\) by \(s+b\) in the characteristic sum. The linear
argument changes from \(\zeta\) to \(\zeta+u\), and the factor
is \(e^{\pi i(b^t\Omega b+2b^t\zeta+a^tb)}\).
Restoring the holomorphic Gaussian this is just \(U_{h^{-1}u}^{(1)}\).
Thus this comparison holds with phases before taking any norm.

For the first sum of the rule write the division points as
\(t=s/N\) with \(s\in\Lambda\otimes\Pi\) modulo \(N\Pi^n\),
as in Lemma~\ref{newtheta:index}. Representatives can be chosen
even in \(\Lambda\otimes\Pi\): its quotient by \(N\Pi^n\) has
odd order since the column is standard at two and \(N\) is odd.
Then \(Ns\in2\Pi^n\) and their semicharacters are trivial.
The finite sum on the initial vacuum is exactly
\(|Z_\Lambda|P_Av\), since \(N\Pi^n\) fixes it with the
prescribed semicharacter. The actions of the second, half-level
sum were computed in Corollary~\ref{bfam:mark}. Equivalently
its translating vectors modulo \(\Lambda\otimes\Pi\) exhaust
\(D/(D\cap A)\): at the primes of \(M\) this is the given
half-level isomorphism, and at its complement the primary torsion
class vanishes. Division representatives may be taken to be zero
modulo \(2N\Pi^n\) at the complement of \(\bar{\mathfrak m}\)
and even in the shifted lattice at \(\bar{\mathfrak m}\), by the
Chinese remainder theorem. Their symmetric lifts consequently
act by the normalized shifts used in that corollary.

In the universal-cover Appell--Humbert frames the constant
biextension frame on two Lie variables is biadditive and
Rosati-compatible. Indeed the factor of automorphy of the
addition biextension is obtained by dividing the factor at
\(z+z'\) by those at \(z,z'\). The resulting cross terms are
additive in the opposite variable, and
\(H(az,z')=H(z,\bar a z')\). Equivalently these statements follow
by taking the rigidified cube trivializations; in this frame the
third differences of the linear and quadratic exponents vanish.
The pullback frame along \((1,w_0)\) is the
Appell--Humbert frame of \(\mathcal L\) just used, by equality
of their Hermitian forms, semicharacters and rigidifications.
Consequently the algebraic jet may be evaluated by ordinary
differentiation on this analytic frame. Any change of compatible
frame is the common bilinear ambiguity killed by harmonicity
in Lemma~\ref{bfam:rule}. Each labelled Hodge derivative is
the corresponding coordinate derivative by the definition of the
framed period chart. Substitution \(z\mapsto Nz\) contributes
\(N^{d_P}\), cancelled by the same denominator in the rule.
All sums are finite except the normally convergent holomorphic
theta series, so these differentiations commute with summation.

The raw column factor is thus \(|Z_\Lambda|/|K_0|\) times the
two square-root indices in Corollary~\ref{bfam:mark}. Setting
\(I=[\Lambda:\Lambda\cap O^n]\), the three factors are
\(I^r,\ |Q|^{-1/2},\ I^{-r/2}\), respectively. This proves
the asserted formula, independent even of the auxiliary \(N\).

For the extension off unit flags it matters that the
Appell--Humbert Gaussian itself need not depend holomorphically
on the period. Work instead in the classical holomorphic
trivialization on the varying torus. The matrices \(P_0,Q_0\)
are holomorphic in the flag coordinate: their entries are either
entries of the inverse flag matrix applied to a fixed period,
or contractions of a flag column against a fixed vector by the
Hermitian form linear in the flag column. Thus \(h,\Omega\) are
holomorphic. The characteristic series without its first Gaussian
defines a holomorphic family of sections on the torus in these
coordinates. Its integral-period automorphy factors, read by
shifting the sum, are holomorphic in \(\Omega,\zeta\); the series
and its derivatives converge normally on compacta, by Gaussian
decay for \(Y\) uniformly positive. Rigidify its line at zero
in this holomorphic trivialization. On each fibre changing to the
Appell--Humbert frame multiplies that function by exactly the
Gaussian in \(T\): it has value one at zero and takes the
classical factors of automorphy to the Appell--Humbert factors.
On the tautological family of period tori this gives the same
relative rigidified polarized line as the pullback of its
biextension by \((1,w_0)\). Indeed their ratio is trivial on
every fibre by the computation of the two Appell--Humbert data;
its
pushforward is a line by cohomology and base change in degree
zero and evaluation is an isomorphism by fibrewise evaluation.
The rigidification trivializes that base line.

Apply the geometric torsion projections and harmonic jet to this
holomorphic characteristic section. These operations are
holomorphic in the period: the torsion points are locally
holomorphic sections of finite étale groups; the formal
biextension frame of Lemma~\ref{bfam:rule} can be constructed
over the local holomorphic base by the same finite-order
operations. The preceding fibre calculation identifies the
result with the analytic marked jet, up to the displayed
nonzero index factors, \(d_L\) and \((\det h)^{n/2}\), all
holomorphic in this chart. Lemma~\ref{dg:gaussian} also proves
this directly by removing the conjugate-type quadratic terms with
the harmonic contractions. This proves holomorphy of the
analytic jet without differentiating a nonholomorphic
\(Y^{-1}\) in the period. Its flag character follows in the
same coordinates. A triangular flag change gives the corresponding
triangular changes in the two labelled tangent blocks (first
columns in one block, last inverse rows in the other). It changes
both period matrices by the same tangent substitution, hence
fixes \(\Omega\); the determinant factor gives the scalar
weight \(n/2\), and the labelled derivatives give the character
of the row-highest polynomial. These are precisely the two parts
of the geometric weight. The geometric reading of a regular
weight rule is holomorphic with that character by the good framed
uniformization (its section property).
Finally all
cross products of the two rows of readings vanish on unit flags
and hence on every flag by the weight and Gram--Schmidt. Where
one analytic reading is nonzero its quotient with the geometric
reading is holomorphic; the quotient has the same absolute value
as on unit flags, hence is constant. Holomorphic continuation
gives the claimed equality everywhere on the component. When
all analytic entries vanish the same reasoning on unit flags
gives zero throughout for the geometric entries as well.
\end{proof}
\src{src/Formalism/LocalToGlobal/ProjectiveLattices/PetersonNormSubspace.lean}{75}{The section property and period identification in a good uniformization};
\src{src/Development/AbelianVarietyTheory/Proofs/SiegelThetaTransformation/PeriodCoordinateThetaOfSignedHeisenbergVacuum.lean}{451}{The half characteristic and determinant in an adapted basis};
\src{src/Development/AbelianVarietyTheory/Proofs/SiegelThetaTransformation/ExistsSymplOfSelfDual.lean}{226}{The probability normalization on the period torus};
\src{src/Development/AbelianVarietyTheory/Proofs/SiegelThetaTransformation/ShortThetaValueComplexPointEqClassicalTheta.lean}{194}{The product of generators and its covariant derivatives};
\src{src/Development/AbelianVarietyTheory/Proofs/SiegelThetaTransformation/ShortSeriesReadingsEqClassicalThetaSeries.lean}{88}{The separated nonzero constants in the actual chart comparison}.

\begin{theorem}[Integral field generation on the marked period charts]\label{bfam:oncharts}
Let \(p,q\ge1,\ m=fr+1,\ n=8m\). Suppose the actual level-\(M\)
PEL cover has at each good embedding a good framed uniformization
by the finitely many marked periods of its common rational
Hermitian space, as in Lemma~\ref{munf:marks} and
Proposition~\ref{bfam:unitread}. Use its identity auxiliary map
and the positive dominant harmonic weight of size \(S=d_P+frn/2\).
Assume the cusp space is read as bounded holomorphic forms of that
weight invariant under the marked principal stabilizers, the raw
rules have the prescribed cusp condition at all good embeddings,
and their readings satisfy the half-maximum hypothesis of
Corollary~\ref{bfam:metric}. Corollary~\ref{crev:good} verifies
these analytic conditions for the selected good covers and their
rule spaces. Use the semiabelian integral tests
of Theorem~\ref{intjet:semi} and the field, level and weight size
bounds of Corollary~\ref{bfam:metric}. Then each algebraic
component has a finite set of integral cusp sections spanning
its cusp space over \(F\). Their logarithmic Petersson norms at
every embedding are at most
\[
                 C(frL)^2 S+2\log D_M .
\]
The same estimate holds for simultaneous generators of the
untwisted full cover with its probability measure.

For the final weight \(\lambda\), put \(W=\sum_{\tau,i}
|\lambda_{\tau,i}|\), \(j_\tau=\lambda_{\tau,p}-1-4m\)
(algebraic weight rows numbered from one, with \(p\) the Lie
multiplicity at \(\tau\)). The first Hermitian block in its period
chart has \(p_0=q\) rows. Use the actual weight decomposition
with \(S\le(1+r)W+rf(1+4m)\), and the key section on the tests of
Theorem~\ref{ctwist:semi}. If \(j\ne0\), assume the horizontal
parallel determinant sections and their product have the common
vertical orders on each local component and the corrected degree
estimate of \eqref{ctwist:degree}, with the stated ramification-one
tests outside \(B\). Their regular trivializations and inverses
with constant pointwise norms preserve the rule cusp conditions
at all good embeddings by Corollary~\ref{crev:twist}. Under the stage bounds
for \(B\), the coefficient field and \(D_M\), the resulting
integral sections span on each component and satisfy
\[
           \log\|f\|_\sigma
            \le C'\bigl((frL)^3 W+(frL)^5\bigr).
\]
If every \(j_\tau=0\), the conclusion holds without the parallel
determinant or degree hypotheses. For \(j\ne0\) the same conclusion
also follows from the one-sided product criterion after
\eqref{ctwist:degree}, provided its lower degree bound gives the
same cost in \eqref{bfam:flat}. On the indefinite principal
family of Proposition~\ref{host:art} with its selected marked
PEL moduli and tests, Theorem~\ref{cmeig:splitmodel} supplies the
common orders and degree for every \(j\). Outside the excluded
primes the original level is hyperspecial, so integrality of
the common order already follows from
Proposition~\ref{cmeig:hyperspecial}.
\end{theorem}
\begin{proof}
First consider the row multiplier. A rational isometry maps the
unit vacuum of a row to a unit scalar times that of the transported
row, by preservation of the specified semicharacter and the
one-dimensional vacuum line. Conjugation also carries the
projectors to those for the transported row, so this scalar is
independent of the column. On the marked stabilizer it is one
at the rank used here. Indeed in a symplectic basis of the
\emph{shifted} row lattice the integral symplectic matrix
preserves its half characteristic, since the isometry commutes
with \(w_0\). For one characteristic theta the residual factor,
after the square root of the Siegel determinant and the quadratic
exponential, is an eighth root of unity. This follows on the
integral symplectic generators. A symmetric integral shear
changes the characteristic by the shear and its half diagonal;
the remaining exponential has argument \(\pi i\) times a
quarter-integer. An integral linear basis change reindexes the
sum, with at most a square root of \(-1\) from its determinant.
Fourier exchange gives \(e^{-\pi i g/4}\) times
\(e^{\pm 2\pi i\alpha^t\eta}\) and the displayed Gaussian
transform in Lemma~\ref{dg:gaussian}'s coordinates.
Integral changes of representatives of the characteristic give
signs, and choices in composing square roots give signs.
Thus the eighth-power assertion persists under composition.
The determinant factors cancel those of the period matrices in
the normalized theta, as in Proposition~\ref{dg:identity}. The
\(n\) identical rows contribute the \(n\)-th power of the
residual factor, equal to one since \(8\mid n\).

The marked translate introduces no extra central character here.
Indeed the residue label is the identity by
Corollary~\ref{bfam:mark}. Both \(y,\gamma y\) lie in
\(2\widehat\Lambda\otimes\bar{\mathfrak m}\widehat L\) for a
stabilizing \(\gamma\), and their residues agree at
\(\mathfrak m\). Their difference lies in \(2D\), by the lift
calculation in Lemma~\ref{bfam:phase}. The vacuum character on
this difference is one. The Weyl phase of the two translations is
also one since their half pairing is integral in the unshifted
self-dual lattice. Hence \(\rho(\gamma y)\) and \(\rho(y)\)
give the very same marked vector. This establishes compatibility
with the invariant geometric readings in the hypotheses. For
arbitrary even \(n\ge8fr+4\) the first conclusion also follows
whenever the geometric and vacuum stabilizer characters match;
divisibility by eight proves this particular matching at the
rank of the final weight composition.

Lemma~\ref{munf:marks} now supplies distinct marked row orbits.
Apply Proposition~\ref{munf:span} with the actual graph identity
of Proposition~\ref{dg:identity}, after dividing each component
by its constant \(C_k\). Equations~\eqref{bfam:reading} and the
nonvanishing of \(u_{\Lambda,N}\) give simultaneous generation.
The sum covers the entire genus of the standard column with its
prescribed level; a single representative need not suffice.
On each chosen column Proposition~\ref{bfam:harm} replaces the
polynomials by its finite spanning set. Lemma~\ref{bfam:rule}
and faithful scalar extension descend generation to \(F\).
Multiplication by the respective reflex scalars and the common
\(D_M^2\) preserves the field span. Theorem~\ref{intjet:semi}
gives integrality, and Corollary~\ref{bfam:metric} gives the
simultaneous norm bounds, including for the component probabilities.

On each component multiply further by its integral key section
of weight \(-4m\) and its integrally rescaled parallel determinant.
Theorem~\ref{ctwist:semi} and equations
\eqref{ctwist:degree}, \eqref{bfam:flat} give the combined
logarithmic cost at most
\[
  2frm\log(2\pi)+12rL\bigl(W+f(1+4m)\bigr)
  +\sum_{\ell\in B}\log\ell+
                    \frac{\log|\Delta_F|}{2[F:\mathbb Q]} .
\]
The last two quantities are \(O((frL)^2)\) under the stage
bounds. Also
\[
 2\log D_M=2(frn+d_P)\log M+2d_P\log\delta_E
                     \ll fL S
\]
by the level and discriminant bounds in
Theorem~\ref{intjet:semi}'s application.
Inserting the stated bound on \(S\) gives the result.
For \(j=0\) the parallel section is identically one and has
neither an order nor a metric cost. If only \(\sum_\tau j_\tau=0\),
formula~\eqref{ctwist:degree} cancels the common degree once
the common finite orders and common corrected degree have been
obtained. Component restriction and extension by zero preserve
the field decompositions and the integral evaluation tests.
\end{proof}
For every selected cover satisfying \(Y.\mathrm{Good}(e,o)\),
Corollary~\ref{crev:good} supplies the marked period, cusp,
half-maximum and measure clauses with the stated stage parameters.
For nonzero parallel weights
Theorem~\ref{cmeig:splitmodel} supplies the local orders and the
arithmetic degree on those marked components. Neither the finite
genus identity nor the positive real reproducing scalar depends
on that degree.
\src{src/Development/AbelianVarietyTheory/Proofs/SiegelThetaTransformation/KernelIdentity.lean}{181}{Doubled identity with the level tests and harmonic insertions};
\src{src/Development/AbelianVarietyTheory/Proofs/SiegelThetaTransformation/HermitianThetaSeriesSpanBoundedForms.lean}{168}{Inversion and reproduction on marked component families};
\src{src/Development/AutomorphicArithmetic/ThetaCorrespondence/ThetaPairLifts/HodgeBundleWeights/ShortSpanCompositionClosed.lean}{69}{The two actual component twists}.

\paragraph{The cusp test on the principal marked charts.}
Retain the period chart and identity Fourier mark of
Propositions~\ref{bfam:unitread} and \ref{bfam:oncharts}.
The following argument uses the mark at one prime over \(M\).
In particular it applies before summing over the definite column genus.

\begin{lemma}\label{crev:mode}
For every rational isotropic row line and every orthonormal frame on
the chart, the real diagonal transvection restriction of each marked
theta jet has mean zero along any period with trivial multiplier.
Its full real cusp-torus average is zero as well.
\end{lemma}
\begin{proof}
Write \(v\ne0\) for the isotropic row and \(I=(Ev)^n\) in the
rational trace symplectic space of the oscillator. Choose a rational
Lagrangian in one row containing \(Ev\), and take its \(n\)-fold sum
\(D\). Complete it to a symplectic polarization \(D\oplus D^*\),
using \(D^*\) for the position variable. Translations by \(a\in
D(\mathbb A_f)\) are multiplication by the finite additive character
of \(z(a)\) at a position \(z\). In particular they act as the identity
at all \(z\) annihilating \(I\).

Let \(\varphi\) be the marked finite vector in these coordinates.
The finite translation calculation preceding Lemma~\ref{bfam:genus}
gives an element of \(I(\mathbb A_f)\) acting on the single marked
vacuum with eigenvalue different from one. Indeed, at a prime of
\(\mathfrak m\) its first \(r\) residues are a basis and a primitive
vector on the conjugate line, divided by the local half-level,
detects one of them. The unshifted character of the isotropic line
is one. It follows pointwise that
\[
                 \varphi(z)=0\qquad (z|_{I(\mathbb A_f)}=0).
\]
This also holds for a sum of the indicated marks, term by term. The
calculation takes place after changing polarization, which intertwines
the Weyl translations by the Fourier and substitution laws in the
proof of Proposition~\ref{bfam:unitread}.

In this polarization the real diagonal flow is a shear. Its
infinitesimal action \(N\) has image \(I_{\mathbb R}\) and kernel
\(I_{\mathbb R}^{\perp}\), and it fixes \(D_{\mathbb R}\) pointwise.
The shear contributes the character with quadratic exponent,
up to the fixed Fourier sign, \(Q(z)=\tfrac12\langle \widetilde
z,N\widetilde z\rangle\), where \(\widetilde z\) is a lift of the
position and the pairing is symplectic. At each place the trace
polarization applied to the transvection
\(\kappa_0 h(y,v)v\) gives a nonzero real multiple of
\(\sum_{j=1}^n|h_\tau(y_j,v)|^2\). The multiples have the same sign:
they are obtained from the product of \(\tau\xi_0\) and
\(\tau\kappa_0=\pm2\tau\xi_0\) with one fixed choice of sign, and
the scaling \(k_0\) of the real parameter is positive. Consequently
\(Q(z)=0\) exactly on the annihilator of \(I_{\mathbb R}\).

More explicitly, on the real flow the first period matrix is fixed
because the shear fixes the chosen first Lagrangian. The second
matrix changes by the first times the real symmetric shear matrix.
Thus \(\Im\Omega\), the holomorphic argument \(hz_j\), the Gaussian
prefactor and \(\det h\) in Proposition~\ref{bfam:unitread} stay
fixed. The change in each summand is exactly its quadratic phase,
also after applying the polynomial in the \(z_j\)-derivatives.
The real restriction therefore has the expansion
\[
            \sum_{\ell\in D^*(\mathbb Q)}
                 \varphi(\ell)\,A_x(\ell)\,e^{2\pi i s Q(\ell)}
\]
with the Fourier sign absorbed in \(Q\). At the fixed frame the
archimedean coefficient is a polynomial times a decaying Gaussian,
and the compact finite support restricts rational positions to a
fixed denominator lattice. The series converges absolutely and
uniformly for real \(s\). If a rational position annihilates \(I\)
at infinity it annihilates \(I\) at all finite places. Hence every
zero-frequency summand vanishes. Averaging over \([0,T]\) and
taking \(T\to\infty\) gives zero by dominated convergence, even
when individual frequencies are not characters of the chosen
period circle. For a periodic sum this is its period mean.
Real translates of the starting frame give the same conclusion.
Integrating along the diagonal circles in the full cusp torus now
gives the torus average by translation invariance of Haar probability.
\end{proof}

\begin{proposition}\label{crev:half}
On each of these principal period charts the raw theta readings at
\(n=8(fr+1)\) are bounded in Petersson norm. For every nonzero
reading each orthonormal frame attaining at least half the component
supremum has row minimum squared at least \(\varepsilon\) of
Proposition~\ref{bergvol:global}. The same minimum conclusion applies
to every nonzero bounded cusp form with trivial transvection multiplier on
the chart. In particular the cusp and half-maximum tests on the raw
rules in Theorem~\ref{bfam:oncharts} follow from its stated marked
period uniformization.
\end{proposition}
\begin{proof}
Write \(a(x)=\min_{0\ne v\in L'}\sum_\tau{\rm maj}_{x,\tau}(v)>0\).
Lemma~\ref{bfam:minimum} and Proposition~\ref{intjet:arch}, at
orthonormalized frames, first give a bound of the form
\[
             \|P(x)\|_{\rm Pet}\le C(1+a(x)^{-1})^B .
\]
Indeed \(\lambda^2\ge 2a(x)/R_E^2\), the theta generator has
mean square one, and the directional norm and its weight correction
are fixed on orthonormal frames. Here \(C,B\) may depend on the
fixed section and stage. Multiplication by its nonzero fixed
reflex scalar does not affect this assertion.

Put \(H=2M\sqrt d\). For \(a=a(x)<\min(1,H^{-1})\) choose a
shortest \(v\). Its integral self-pairing vanishes by the norm
argument following Proposition~\ref{bergvol:global}. The
transvection there acts trivially modulo \(M\) on \(\mathfrak a L'\),
so lies in the actual marked stabilizer by Lemma~\ref{munf:marks}.
Its multiplier on the raw sections is trivial by the eighth-power
and identity-mark calculation in Theorem~\ref{bfam:oncharts},
which uses only the transformation formula. Take \(H\) as period
in the scaled coordinate, and write \(G(\zeta)\) for the raw
coefficient along this flow. It is holomorphic on
\(\Im\zeta>-1/a\) by the leading-minor calculation above.

We give the majorant comparison needed to apply the polynomial bound
before knowing a global supremum. At a single place set
\(l=\|\tau v\|_x^2\le a\). Rotate the positive and negative blocks
so the two components of \(\tau v\) lie on their first coordinate
axes, each of norm \(\sqrt{l/2}\). These rotations preserve the
central majorant. In this plane the new positive line has positive
and negative coordinates \(A=1-i\zeta l/2,\ B_1=-i\zeta l/2\),
with \(|A|^2-|B_1|^2=1+l\Im\zeta\). All orthogonal directions
retain their majorant. The two comparison eigenvalues in the plane
are \(K\) and \(K^{-1}\), where
\[
 K=\frac{(|A|+|B_1|)^2}{1+l\Im\zeta}
    \le\frac{(1+l|\zeta|)^2}{1+l\Im\zeta}.
\]
Indeed a unitary change of phase makes the normalized line coordinates
nonnegative and its orthogonal negative mate interchanges them.
At \(\zeta=s-i/(2a)\), \(0\le s\le H\), this gives \(K\le 13\):
the denominator is at least \(1/2\) and
\(l|\zeta|\le Ha+1/2\le3/2\).
Thus the new row minimum is at least \(a/13\).
The leading-minor variation estimate preceding
Lemma~\ref{bfam:harm} bounds the conversion from Petersson norm
to raw coefficient on this line by \(2^{w/2}\).
It follows that
\[
      \sup_{0\le s\le H}|G(s-i/(2a))|
                    \le 2^{w/2}C(1+13/a)^B .
\]
For \(y=\Im\zeta\ge0\), \(0\le s\le H\), the same comparison
gives \(K\le(1+a(H+y))^2\). The leading-minor formula bounds
the raw metric conversion by \((1+ay)^{w/2}\).
Together these give polynomial growth of \(G(s+iy)\) in \(y\)
for the fixed starting frame. In \(q=\exp(2\pi i\zeta/H)\),
Cauchy's coefficient estimates therefore remove all negative
Laurent powers. Lemma~\ref{crev:mode} removes the constant power.
The maximum principle for \(G/q\) on the disk of radius
\(\exp(\pi/(aH))\) gives
\[
      \|P(x)\|_{\rm Pet}=|G(0)|
       \le 2^{w/2}C(1+13/a)^B e^{-\pi/(aH)} .
\]
This is bounded as \(a\downarrow0\). On the complementary range
the original polynomial estimate is bounded. Hence every nonzero
section has a finite positive supremum and one can choose an
orthonormal frame above half of it. With boundedness established,
the half-maximum argument following Proposition~\ref{bergvol:global}
applies with its original \(\varepsilon\); the qualitative constants
\(B,C\) disappear. For a general bounded cusp form the real torus
and positive-frequency argument there supplies its diagonal zero
mode. Thus the same quantitative conclusion holds for that space
at every weight whose transvection character is trivial.
\end{proof}
The proof uses the short vector of the component under consideration.
It gives the absolute-constant bound of Corollary~\ref{bfam:metric}
for all selected parameters, all permitted column lattices and
auxiliary integers at every good embedding; it requires neither
a common fundamental domain as the rank changes nor a uniform
constant in the preliminary polynomial growth estimate.
\src{src/Development/AbelianVarietyTheory/Proofs/CuspFormEvaluations/ExistsExpansionOfEqSumThFr.lean}{26}{The source's different zero-mode test using positivity of the last row weight};
\src{src/Development/AbelianVarietyTheory/Proofs/DivisorProofs/FirstMinAtHalfMax.lean}{111}{Transfer to the minimum at the selected level}.

\begin{corollary}\label{crev:twist}
On the good marked charts the algebraic rule space defined by
testing all good complex readings requires no additional cusp
comparison for the raw spanning argument. A regular unit section on
a component with constant positive Petersson norm, in mutually dual
multiplicative metrics for it and its inverse, preserves this cusp
condition in both directions. This applies to the key and the
horizontal determinant whenever they exist on that component with
their stated constant metrics.
\end{corollary}
\begin{proof}
In the rule definition the cusp subspace is the intersection of
the regular weight rules with the inverse images, under every good
reading, of the complex spans of the analytic cusp forms. Each such
span is the analytic cusp space itself: holomorphy, transformation
law and vanishing of the torus integrals are linear, and each
finite sum of bounded forms is bounded. Conversely inclusion in the
span holds for each member. Lemma~\ref{bfam:rule},
Lemma~\ref{crev:mode} and Proposition~\ref{crev:half} place the
raw rules in this intersection. For generation one only needs to
express the reading of every member in the span of the raw readings,
which is the conclusion of Proposition~\ref{munf:span} on those
charts. Injectivity after scalar extension in Lemma~\ref{bfam:rule}
then descends the span to the coefficient field.

For the multiplication assertion use a common finite-index
transvection period lattice killing both characters. Passing to
it leaves each normalized constant term unchanged when applied to
a function periodic for the original lattice. Along a rational
isotropic row the positive-ray metric estimate following
Proposition~\ref{bergvol:global} and boundedness imply that
every nonzero Fourier coefficient of either factor has all
coordinates nonnegative, by growing one imaginary coordinate
at a time. At height zero the coefficients are absolutely summable
by smoothness on the real torus. Thus the product constant is the
product of the two constants, by convolution: two vectors with
nonnegative coordinates can sum to zero only if both are zero.
This proves cusp vanishing for the product, and boundedness follows
by multiplicativity of the metric. Apply the same argument to the
inverse section, which is holomorphic and has the reciprocal
constant norm. In particular its Fourier support is tested afresh
using holomorphy in the same upper tube, even though its \emph{weight}
has the opposite sign. The polynomial growth bound uses only the
weight spread. This treats the Fourier condition in the rule space;
a separate use of a subcanonical sheaf on a compactification must
identify its restriction with that space.
\end{proof}
\src{src/Formalism/LocalToGlobal/ProjectiveLattices/CuspFree.lean}{14}{All good embeddings in the definition of the cusp rules};
\src{src/Formalism/LocalToGlobal/ProjectiveLattices/JFDFree.lean}{42}{The span of the bounded holomorphic cusp readings}.

\paragraph{The flag measure of the selected unitary genus.}
Let \(F_0=E^+\). Use the host fields of Proposition~\ref{host:art},
the genus of the self-dual row, and its odd principal level \(M\ge3\).
Write \(K\) for the finite adelic stabilizer and \(K(M)\) for its
principal subgroup. Let \(c_r\) and \({\rm mass}_I\) be as in
Proposition~\ref{bergvol:compactmass}.

\begin{proposition}\label{crev:mass}
For the full adelic packet of marked rows, the sum of the component
flag volumes in the normalization of Lemma~\ref{bergvol:local} is
\[
        V=[K:K(M)]\,\frac{c_r^f\,{\rm mass}_I}{(2\pi)^{fr}}
           \le V_* .
\]
Adelic probability on flag-invariant functions pushes forward to
the sum of the fundamental-set measures divided by \(V\).
Conditioning on a nonempty union of components gives the same
formula with its own total flag volume.
\end{proposition}
\begin{proof}
In the marked period description of Lemma~\ref{munf:marks}, the
finite double cosets of \(G=U(h_0)\) represent the marked rows of
this genus. Indeed finite translates of the completed reference
lattice and its basis give a row by intersection with the rational
space. Conversely choose local isometries to a specified row in
the genus. At primes of \(M\) the extension is split and the mark
is a basis on one half with its pairing-dual on the other. An
invertible matrix over the local residue quotient lifts to an
integral invertible matrix by lifting entries. Thus the local
isometries may match the marks and the stabilizer is \(K(M)\).
In terms of the two algebraic Lie ranks \(p,q\), the real flag
space at a place is \(U(p_0,r-p_0)/U(1)^r\), with \(p_0=q\).

Use the rational form \(J=\operatorname{diag}(1_{p_0},-1_{r-p_0})\) given
by Lemma~\ref{cmeig:row}, and compare with \(G_c=U(1_r)\), the
group of the positive definite standard form. At every finite
place the two self-dual lattices are isometric, including residue
characteristic two, by the local part of that lemma. Keep track of
the Haar densities in this assertion. On the Lie algebra for a
Hermitian matrix \(J\), left translation identifies an invariant
top differential by identifying its tangent at one with the fixed
\(F_0\)-space of skew-Hermitian matrices via \(X\mapsto JX\).
Choose the same rational top differential on that vector space
for both \(J\) and \(1_r\). Under an isometry \(P^*JP=1_r\),
the comparison in these coordinates is \(A\mapsto P^*AP\).
Its determinant on the skew-Hermitian space is
\(N_{E/F_0}(\det P)^r\): extend scalars to the split algebra,
where the two factors multiply an unrestricted \(r\)-by-\(r\)
matrix on its left and right. Its absolute value is one at each
finite place because \(N(\det P)=1/\det J=\pm1\).
Thus the corresponding maximal stabilizers have the same local
volume. No choice of their local bases contributes a discriminant.

The global scalar and convergence factors in the Tamagawa measure
are the same: the base field and dimension agree and the geometric
character lattice in both cases is generated by determinant with
the quadratic Galois action of \(E/F_0\). At each real place the
common gauge in the \(A\)-coordinates is the same scalar times
Lebesgue measure on the imaginary diagonal and the complex
strictly lower entries. Multiplication by \(J\) only changes signs
in these coordinates. Divide out that same scalar for both groups.
Use volume one for the maximal stabilizers at finite places and
the resulting real coordinate densities. In these choices the
total adelic covolumes of the two groups are equal: they are
obtained from Tamagawa measure by a common scalar and both
Tamagawa numbers are two for a nonzero Hermitian space over a
quadratic extension. See \href{https://doi.org/10.1007/BF02391774}
{Weil, \emph{Sur la formule de Siegel dans la th\'eorie des groupes
classiques}}, and Ichino, \emph{A regularized Siegel--Weil formula
for unitary groups}, \S4. For the definite group the real volume is
\(c_r^f\). The finite double cosets at maximal level are its
standard lattice genus: cutting the quotient over such a coset
divides this common real volume by its finite automorphism group.
Summing gives \(c_r^f\,{\rm mass}_I\).

For the indefinite group, the quotient of the real density by
angle measure \(d\theta_1\cdots d\theta_r\) is the stated flag
measure. At the central flag the tangent projection takes the
strictly lower entries of the Lie algebra with Jacobian one.
At a lower unipotent matrix \(u=k b\), with \(k\) unitary and
\(b\) upper triangular, the projected tangent in this frame is
the strictly lower part of \(b(u^{-1}du)b^{-1}\).
The unipotent differential has determinant one in its entries.
Conjugation on matrices modulo upper triangular matrices has
complex determinant \(\prod_{i>j} b_{ii}/b_{jj}\). Its squared
absolute value is \(\prod_{i=1}^{r-1}D_i^{-2}\), using
\(|b_{ii}|^2=D_i/D_{i-1}\) and \(D_0=D_r=1\).
This recovers the chart density at every point.

Over each finite double coset \(gK(M)\), the quotient of
\(G_\infty\times gK(M)\) is by
\(\Gamma_g=G(F_0)\cap gK(M)g^{-1}\). Cutting the real factor
by a fundamental set for \(\Gamma_g\) gives the common finite
volume \(1/[K:K(M)]\) times real quotient volume. There is no
stabilizer on the flag: it would be discrete in a compact torus,
hence finite, and the principal subgroup on the underlying
integral period lattice is torsion-free. For the latter, at an
odd prime dividing \(M\) the powers of a nonidentity
\(1+\ell^a B\), \(a\ge1,\ B\not\equiv0\bmod\ell\), have depth
\(a+v_\ell(t)\) at exponent \(t>0\), by the binomial theorem.
Integrating along the angle tori therefore gives \((2\pi)^{fr}\)
times the flag integral. The same disintegration applies to any
nonnegative measurable function of the flag. Applying it first
to one proves the identity and then dividing by the total mass
proves the probability formula. The congruence index is at most
\(M^{2fr^2}\) by reduction on the row, so
Proposition~\ref{bergvol:compactmass} gives the required bound.
\end{proof}
In particular both analytic tests in
Proposition~\ref{bergvol:global} are available on the selected
principal packet for the bounded analytic cusp space of trivial
transvection character. Using this for an algebraic component
uses its identified framed period charts. For the rule space the
reading is as in Corollary~\ref{crev:twist}; agreement with a
separately defined toroidal sheaf is a geometric comparison.

\begin{corollary}[Analytic estimates on a selected good cover]\label{crev:good}
Take a stage in the selected range on the host of
Proposition~\ref{host:art}, a core with the prescribed host,
polarized datum and level, and the deep cover \(Y\) with
\(Y.\mathrm{Good}(e,o)\). Use its identity auxiliary rule
at \(n=8(fr+1)\), and the parameters and semiabelian tests of
Theorem~\ref{bfam:oncharts}. At every good embedding the raw
rules are cusp rules on this \(Y\). The short integral rules in
that theorem span the untwisted space on the cover and on each
component, with logarithmic Petersson norms at most
\(C(frL)^2S+2\log D_M\) for the stated stage size bounds.
For the component cusp spaces at the frame weights of
Corollary~\ref{bergvol:size} the normalized evaluation kernel
is at most \(\beta_*\).
\end{corollary}
\begin{proof}
Here \(Y.\mathrm{Good}(e,o)\) contains the following data. Its
\code{unf} gives the complex framed period uniformization at each
good embedding, including the section transformation law and the
Haar mean-square-one normalization of the theta generator on unit
frames. In these coordinates the first \(p_0=q\) Hermitian
directions are positive. Their conjugate, reversed tangent block
has Lie rank \(q\), and the other tangent block has Lie rank \(p\)
at the labelled \(\tau\). Thus the pivot \(\lambda_{\tau,p}\)
in Theorem~\ref{bfam:oncharts} is numbered in the algebraic
weights, while the positive block used for the majorant ends at
\(p_0\) in the Hermitian flag.

The \code{modelW} and \code{xiAbs} clauses give the signed host
Gram matrix and the absolute polarization coefficient at every
place. The common lattice embedding and alternating pairing in
\code{adPin} identify the original period lattice, its finite
adelic translates, the half-level mark and the eigenframe
coordinates on these same components. In particular stripping the
fixed host ideal from a translate of the reference period
\(\mathfrak a L'_*\) gives the translated Hermitian self-dual row:
finite unitary matrices commute with \(\mathfrak a\), and
localization detects self-duality. This is \(L'\) in the return
count and in Lemma~\ref{bfam:minimum}. The \code{lev} clause gives
auxiliary integer one on the deep variety; the spanning construction
there uses its identity auxiliary map of Corollary~\ref{bfam:mark}.
It does not replace the original map from the deep variety to the
base by an identity.

The \code{dictU} and \code{levelK} clauses identify the rational
and right-principal-level equivalences with the chart action and
exhaust the real flag orbits. Indeed the adelic frame map is
equivariant modulo arithmetic translates under the rational and
level actions and intertwines the real action; it hits each flag
modulo its upper triangular frame change. Equality of its frames
modulo the arithmetic action forces rational and finite-level
equivalence. The unit frames over each flag form a single orbit
under its diagonal angle group. Together with the marked-period
classification in Lemma~\ref{munf:marks}, this is the packet and
the probability disintegration of Proposition~\ref{crev:mass}.
The \code{unit} clause gives leading Gram determinants of absolute
value one at the adelic frames, as used in their metric reading.

Finally \code{chiTriv} gives trivial arithmetic multiplier at the
dominant natural weights on each good chart. It gives the same
for every integral flag weight: for fixed arithmetic element its
multiplier is multiplicative in addition of weights and never
zero. Take natural weights with \(n=0\) that are one on an
initial segment at one labelled embedding and zero elsewhere.
Consecutive differences give all individual weight coordinates,
with reversed order and a minus sign on the dual block. Their
integral span is the whole flag weight lattice. In particular
there is no character loss in the transvection period for the
frame weights. Apply Propositions~\ref{crev:half} and
\ref{crev:mass}, then Proposition~\ref{bergvol:global}.
The actual all-weight space \code{cuspW} is the intersection of
the regular rules with the good-reading tests of
Corollary~\ref{crev:twist}; at a natural weight it is the ordinary
rule space \code{cusp}, and \code{cuspOn} adds only the component
support condition. Thus Theorem~\ref{bfam:oncharts} applies with
its integral tests and no additional analytic cusp assumption.
\end{proof}
\src{src/Development/AbelianVarietyTheory/Definitions/Divisors/NormNegCmConstantEqAtAllEmbeddings.lean}{187}{The clauses of the selected deep cover};
\src{src/Formalism/LocalToGlobal/ProjectiveLattices/FamilyUniformisationBasePointOrthogonalBis.lean}{97}{The marked finite translates in the period identification};
\src{src/Development/AbelianVarietyTheory/Definitions/Foundations/IsProbabilityMeasureMuPart.lean}{425}{The actual all-weight rule space};
\src{src/Development/AbelianVarietyTheory/Counterexamples/LevelGroupExamples/AdelEqZeroOfOffDom.lean}{77}{Generation of the arithmetic characters by natural weights}.

Constructing such covers at all scales is a separate existence
statement. For example, the sized-cover construction starts with
the pinned core, a stage in range, the host splitting bound at least
five and \(15\) coprime to the base level. It constructs the half-level,
uses PEL moduli representability, forms the coefficient compositum
with a ray class field, and builds the marked uniformization; the
field-size bound uses the ray-class-field estimate. If the comparison
to a prescribed base frame requires \(Y.\mathrm{Match}\), that
condition concerns the same base embedding and base frames and is
not part of \(\mathrm{Good}\). The corollary asserts the estimates
on the selected \(Y\); it does not construct the core, the compatible
family or its toroidal compactifications.
\src{src/Development/AbelianVarietyTheory/Proofs/DivisorProofs/ExistsGoodCoverOverSizedField.lean}{111}{Inputs and construction of the sized good level};
\src{src/Development/AbelianVarietyTheory/Definitions/Divisors/NormNegCmConstantEqAtAllEmbeddings.lean}{207}{The common marked frames included in goodness}.

